\section{Zaključak}

Distribuirane glavne knjige zasnovane na usmerenom acikličnom grafu predstavljaju značajnu inovaciju u DLT oblasti, nudeći alternativu tradicionalnim blokčejn arhitekturama. Prednosti koje ovi sistemi donose jasno ukazuju na njen potencijal, posebno u aplikacijama koje zahtevaju obradu velikih količina podataka u realnom vremenu i mikrotransakcije.

Međutim, značajni izazovi i ograničenja otežavaju primenu ove tehnologije van specifičnih i eksperimentalnih okruženja. Ovo ukazuje na to da je tehnologija još uvek u fazi razvoja i testiranja pre nego što postane univerzalno prihvaćena.

Dalji razvoj standardizovanih protokola, unapređenje bezbednosnih karakteristika i eksperimentalne primene u realnim slučajevima upotrebe mogli bi omogućiti da, u određenim domenima, DAG strukture postanu konkurentan izbor tradicionalnim blokčejn sistemima.